% Independently written reductions for integration into the GRH paper.
% Required packages: amsmath, amssymb, amsthm.
% Required theorem environments: theorem, lemma, proposition, corollary,
% definition, remark. Bibliography keys: OpenAIOctober, OpenAISeptember.

\section{Smoothed sums and zeros of L-functions}
\label{qr-sec-mellin}

The final analytic step can be separated from the arithmetic estimates.
This is useful because a proposed simplification must supply a precisely
stated estimate; an analogy with a known proof does not supply that estimate.
We first prove the implication that we shall need. The argument is related
to the reduction in \cite[Section~3]{OpenAIOctober}.
% Source labels: sec:reduction, eq:mobius-saving.

Let $K$ be a number field, let $\mathcal O_K$ be its ring of integers,
and let $N\mathfrak a=|\mathcal O_K/\mathfrak a|$ denote the norm of a
nonzero integral ideal. Here a number field is a finite extension of $\mathbb Q$, and
$\mathcal O_K$ consists of its elements integral over $\mathbb Z$.
For clarity, a ray class group is the group of fractional ideals
coprime to a fixed nonzero ideal $\mathfrak f$, modulo the principal
ideals $(\alpha)$ for which $\alpha\equiv1\pmod{\mathfrak f}$
locally at each prime dividing $\mathfrak f$, with positivity imposed
at a chosen set of real embeddings. Locally the congruence means
$v_{\mathfrak p}(\alpha-1)\geq v_{\mathfrak p}(\mathfrak f)$,
where $v_{\mathfrak p}$ is the exponent of $\mathfrak p$ in a
principal fractional ideal.
A finite-order Hecke character $\nu$ is a homomorphism of such a
ray class group into the complex roots of unity. We view it as a
completely multiplicative function on integral ideals coprime to
$\mathfrak f$ and extend it by zero to all other ideals. The least
possible finite modulus is its conductor, and we use that modulus.
In particular, $|\nu(\mathfrak a)|\leq1$. The ideal M\"obius function is
defined by $\mu_K(\mathcal O_K)=1$, by $\mu_K(\mathfrak a)=0$ if a square
of a prime ideal divides $\mathfrak a$, and by
$\mu_K(\mathfrak p_1\cdots\mathfrak p_r)=(-1)^r$ for distinct prime ideals.
We use the classical analytic properties of the associated $L$-function:
its Euler product converges absolutely in $\Re s>1$, and it has a
meromorphic continuation. The result below is an implication from these
properties and a stated sum estimate; it does not prove that estimate.

For a finite set $S$ of prime ideals, write $(\mathfrak a,S)=1$ if none
of them divides $\mathfrak a$, and put
\begin{equation}\label{qr-eq-deleted-L}
 L_K^S(s,\nu)=L_K(s,\nu)
   \prod_{\mathfrak p\in S}
       \bigl(1-\nu(\mathfrak p)(N\mathfrak p)^{-s}\bigr).
\end{equation}
Every factor in this finite product is nonzero in $\Re s>0$.
For $W\in C_c^\infty((0,\infty);\mathbb C)$ and $X>0$, define
\begin{equation}\label{qr-eq-smooth-sum}
 A_W(X)=\sum_{(\mathfrak a,S)=1}
       \mu_K(\mathfrak a)\nu(\mathfrak a)
       W(N\mathfrak a/X).
\end{equation}
The sum is finite, since the norm is bounded on its support.

\begin{theorem}[A smoothed cancellation criterion]\label{qr-thm-mellin}
Let $\theta\geq0$. Suppose that for every $\varepsilon>0$ and every
$W\in C_c^\infty((0,\infty);\mathbb C)$ there is a constant
$C_{\nu,S,W,\varepsilon}$ such that
\begin{equation}\label{qr-eq-cancellation-input}
 |A_W(X)|\leq C_{\nu,S,W,\varepsilon}X^{\theta+\varepsilon}
 \qquad(X\geq2).
\end{equation}
Then $L_K(s,\nu)$ has no zero in $\Re s>\theta$.
\end{theorem}

\begin{proof}
Fix a test $W$ supported in $[a,b]\subset(0,\infty)$. Since a nonzero
integral ideal has norm at least one, $A_W(X)=0$ for $X<1/b$. Define
\[
 \mathcal M_W(s)=\int_0^\infty A_W(X)X^{-s}\,\frac{dX}{X}.
\]
This integral is locally uniformly convergent in $\Re s>\theta$.
Indeed, for a compact subset of that half-plane choose $\delta>0$ such
that $\Re s\geq\theta+2\delta$ there, and apply
\eqref{qr-eq-cancellation-input} with $\varepsilon=\delta$.
The integrand at infinity is then bounded by a constant times
$X^{-\delta}dX/X$. Derivatives with respect to $s$ introduce fixed
powers of $\log X$ and remain integrable. The interval near zero makes
no contribution. Thus $\mathcal M_W$ is holomorphic in $\Re s>\theta$.

The Mellin transform
\[
 \widehat W(s)=\int_0^\infty W(y)y^s\,\frac{dy}{y}
\]
is entire. For $\Re s>\max(1,\theta)$, absolute convergence of the
Euler product permits interchange of the ideal sum and the integral.
More explicitly, the sum of the integrals of absolute values is
\[
 \left(\int_0^\infty |W(y)|y^{\Re s}\,\frac{dy}{y}\right)
 \sum_{(\mathfrak a,S)=1}
 \frac{|\mu_K(\mathfrak a)\nu(\mathfrak a)|}{(N\mathfrak a)^{\Re s}},
\]
which is finite. The substitution $y=N\mathfrak a/X$ therefore gives
\begin{equation}\label{qr-eq-mellin-identity}
 \mathcal M_W(s)=\widehat W(s)
       \sum_{(\mathfrak a,S)=1}
       \frac{\mu_K(\mathfrak a)\nu(\mathfrak a)}{(N\mathfrak a)^s}
       =\frac{\widehat W(s)}{L_K^S(s,\nu)}.
\end{equation}
The last equality follows by expanding the absolutely convergent
reciprocal Euler product. The identity theorem for meromorphic functions
now gives
\begin{equation}\label{qr-eq-continued-product}
 L_K^S(s,\nu)\mathcal M_W(s)=\widehat W(s)
 \qquad(\Re s>\theta).
\end{equation}
At a pole this is understood as an identity of meromorphic functions.

Suppose that $L_K(\rho,\nu)=0$ and $\Re\rho>\theta$. Since
$\theta\geq0$, the factors deleted in \eqref{qr-eq-deleted-L} are
nonzero at $\rho$, so $L_K^S(\rho,\nu)=0$. Choose a nonzero,
nonnegative function $\phi\in C_c^\infty((1,2);\mathbb R)$ and set
$W(y)=\exp(-\rho\log y)\phi(y)$, where $\log y$ is the real logarithm.
This is an admissible test, and
\[
 \widehat W(\rho)=\int_1^2\phi(y)\,\frac{dy}{y}>0.
\]
But \eqref{qr-eq-continued-product}, evaluated at $\rho$, says that this
positive number is zero. This is a contradiction.
\end{proof}

\begin{remark}\label{qr-rem-mellin-quantifiers}
The constants in \eqref{qr-eq-cancellation-input} may depend on the
fixed character and on the test. The test in the proof is chosen only
after the hypothetical zero is fixed. The theorem excludes an open
half-plane and makes no assertion at its boundary. It does not require
an unsmoothed M\"obius estimate or a bound for $1/L_K$ near a zero.
\end{remark}

For $K=\mathbb Q$, the same proof uses positive integers, a Dirichlet
character $\chi$, and the usual M\"obius function. Thus the bound
\eqref{qr-eq-cancellation-input} with $\theta=1/2$, for every primitive
Dirichlet character and every fixed test, implies GRH by
Theorem~\ref{prim:strip-reduction}. Its functional equation and the
imprimitive Euler factors have already been proved there. The new
arithmetic task would be the square-root estimate itself.

\section{Extracting one sum from a second moment}
\label{qr-sec-moments}

The following elementary lemma makes explicit how the size of a family
affects the exponent obtained for one of its sums.

\begin{lemma}\label{qr-lem-replicated-rows}
Let $(A_u)_{u\in\mathcal R}$ be a finite family of complex numbers,
and let $A$ be another complex number. Suppose there are $J\geq1$
distinct indices $u_1,\ldots,u_J$ such that $|A_{u_j}-A|\leq E$,
where $E\geq0$. Then
\begin{equation}\label{qr-eq-replicated-rows}
 |A|^2\leq\frac2J\sum_{u\in\mathcal R}|A_u|^2+2E^2.
\end{equation}
\end{lemma}

\begin{proof}
For each $j$, the triangle inequality and $(x+y)^2\leq2x^2+2y^2$
give $|A|^2\leq2|A_{u_j}|^2+2E^2$. Average over $j$ and enlarge
the sum of nonnegative terms to all of $\mathcal R$.
\end{proof}

\begin{proposition}\label{qr-prop-power-rows}
Fix an integer $m\geq2$. Suppose, at parameters $D\geq2$ and $H>1$,
that a family satisfies
\[
 \sum_{u\in\mathcal R(H)}|A_u(D)|^2\leq C_\eta D^{1+\eta}H
\]
for an arbitrarily prescribed $\eta>0$. Put $Y=H^{1/m}$ and suppose
that $Y\geq2$. Assume there are at least $cY/\log Y$ distinct indices
in this family at which $A_u(D)$ differs from a distinguished sum
$A_1(D)$ by at most $C D/Y$. Then
\begin{equation}\label{qr-eq-power-extraction}
 |A_1(D)|^2\ll_\eta
 D^{1+\eta}H^{1-1/m}\log H+D^2H^{-2/m}.
\end{equation}
The implied constant may depend on $m,c,C,C_\eta$.
\end{proposition}

\begin{proof}
Apply Lemma~\ref{qr-lem-replicated-rows} with
$J\geq cH^{1/m}/\log(H^{1/m})$ and $E\leq CDH^{-1/m}$.
\end{proof}

\subsection{The sextic example}
\label{qr-subsec-sextic}

We state the arithmetic hypotheses carefully to separate a proved
deduction from its analytic input. Put
$K=\mathbb Q(\sqrt{-3})$, $\omega=e^{2\pi i/3}$, and
$\mathcal O=\mathbb Z[\omega]$. For $z=a+b\omega$ its norm is
$N(z)=|z|^2=a^2-ab+b^2$.
The ring $\mathcal O$ is Euclidean for this norm: writing a complex
number as $x+y\omega$ and rounding $x,y$ to the nearest integers gives
an error $u+v\omega$ with $|u|,|v|\leq1/2$, whose norm is at most
$3/4<1$. Applied to a quotient of two elements, this gives division
with remainder of smaller norm. Hence every ideal is principal.
The six units are $\{\pm1,\pm\omega,\pm\omega^2\}$; their images are
the six distinct units of $\mathcal O/(3)$. Consequently every ideal
coprime to $3$ has a unique generator congruent to $1$ modulo $3$.
We call this its primary generator.

For a prime ideal $\mathfrak p$ not dividing $6$ and an element $u$
not in $\mathfrak p$, the sextic residue symbol
$(u/\mathfrak p)_6$ is the sixth root of unity whose reduction modulo
$\mathfrak p$ equals $u^{(N\mathfrak p-1)/6}$. Set it equal to zero
when $u\in\mathfrak p$, and extend multiplicatively in the ideal
denominator. For the unit ideal the symbol is one for every $u$.
These definitions are possible because the six roots of unity remain
distinct modulo $\mathfrak p$; in particular $6$ divides
$N\mathfrak p-1$.

Fix a finite-order Hecke character $\nu$ and a finite set $S$ containing
the primes above $2$ and $3$ and those dividing its conductor. Define
\begin{equation}\label{qr-eq-sextic-family}
 A_u(D)=\sum_{(\mathfrak n,S)=1}
     \mu_K(\mathfrak n)\nu(\mathfrak n)
     (u/\mathfrak n)_6 W(N\mathfrak n/D).
\end{equation}
The sum $A_1(D)$ is the sum in \eqref{qr-eq-smooth-sum}.

\begin{proposition}[Conditional sextic estimate]\label{qr-prop-sextic}
Assume the following two inputs for this family.
First, for each sufficiently large $Y$, at least $c_S Y/\log Y$
prime ideals outside $S$ satisfy $Y/2<N\mathfrak p\leq Y$, where
$c_S>0$ is fixed. Second, for every fixed $0<\vartheta\leq1/10$,
every $\eta>0$, and every fixed test $W$, one has
\begin{equation}\label{qr-eq-sextic-ms-input}
 \sum_{0<N(u)\leq D^{1+\vartheta}}|A_u(D)|^2
 \ll_{\nu,S,W,\vartheta,\eta}D^{2+\vartheta+\eta}
 \qquad(D\geq2).
\end{equation}
Then for every $\varepsilon>0$ and every fixed test $W$,
\begin{equation}\label{qr-eq-sextic-conclusion}
 A_1(D)\ll_{\nu,S,W,\varepsilon}D^{11/12+\varepsilon}.
\end{equation}
Consequently $L_K(s,\nu)$ has no zero in $\Re s>11/12$.
\end{proposition}

The prime-counting input is the stated consequence of the prime ideal
theorem in this fixed field. Estimate \eqref{qr-eq-sextic-ms-input}
is the principal input asserted in \cite[Section~3]{OpenAIOctober};
that source also states explicit control by a finite number of
derivatives of $W$. Its automorphic proof is not reproduced here.
The following proof establishes the implication from the displayed
inputs, independently of that automorphic proof.
% Source labels: thm:ms, eq:ms, eq:prime-extract.

\begin{proof}
For a prime ideal $\mathfrak p$ outside $S$, let $p$ be its primary
generator. The definition of the residue symbol gives exactly
\[
 (p^6/\mathfrak n)_6=\mathbf1_{\mathfrak p\nmid\mathfrak n}.
\]
Therefore the difference $A_1(D)-A_{p^6}(D)$ contains only ideals
divisible by $\mathfrak p$.

For completeness, the required counting bound is elementary in this
field. The inequality
$a^2-ab+b^2\geq(a^2+b^2)/2$ shows that the number of elements of
$\mathcal O$ of norm at most $T$ is $O(T)$ for $T\geq1$. Since every
ideal is principal, the same upper bound holds for ideals. If $W$ is
supported in $[a,b]$, multiplication by $\mathfrak p$ identifies the
ideals in question with ideals of norm at most $bD/N\mathfrak p$.
If this last quantity is below one there are none; otherwise counting
and $|\mu_K\nu|\leq1$ give, for $Y/2<N\mathfrak p\leq Y$,
\begin{equation}\label{qr-eq-sextic-error}
 |A_1(D)-A_{p^6}(D)|\ll_W D/Y.
\end{equation}

Put $H=D^{1+\vartheta}$ and $Y=H^{1/6}$. The selected elements
$p^6$ have norm at most $H$ and are distinct. Indeed, equality of
sixth powers implies equality of the principal ideals $(p)^6$,
and unique factorization of ideals then implies equality of the
prime ideals. By the assumed prime count there are
$J\gg_S Y/\log Y$ such rows for large $D$.
Lemma~\ref{qr-lem-replicated-rows} gives
\begin{align}
 |A_1(D)|^2
 &\ll_{\nu,S,W,\vartheta,\eta}
 D^{2+\vartheta+\eta}Y^{-1}\log Y+D^2Y^{-2}\notag\\
 &\ll_{\nu,S,W,\vartheta,\eta}
 D^{11/6+5\vartheta/6+\eta}\log D
       +D^{5/3-\vartheta/3}.
 \label{qr-eq-sextic-squared}
\end{align}
For a given $\varepsilon>0$, choose
$\vartheta=\min(1/10,\varepsilon/5)$ and $\eta=\varepsilon/2$.
Use $\log D\ll_\varepsilon D^{\varepsilon/2}$ and take square roots.
The first exponent is at most
$11/12+\varepsilon/12+\varepsilon/2<11/12+\varepsilon$;
the second is at most $5/6$. This proves
\eqref{qr-eq-sextic-conclusion} for large $D$. The bounded remaining
range follows from ideal counting, after increasing the constant.
Theorem~\ref{qr-thm-mellin} gives the zero-free conclusion.
\end{proof}

\subsection{A limitation of this second-moment argument}
\label{qr-subsec-exponent-limit}

\begin{proposition}\label{qr-prop-exponent-limit}
For $m\geq2$ and $\alpha>0$, put
\[
 E_m(\alpha)=\max\left\{
     \frac12+\frac{m-1}{2m}\alpha,\quad
     1-\frac\alpha m\right\}.
\]
Then
\begin{equation}\label{qr-eq-optimal-exponent}
 \min_{\alpha>0}E_m(\alpha)=\frac{m}{m+1},
 \qquad\alpha_{\min}=\frac{m}{m+1}.
\end{equation}
In the restricted range $\alpha\geq1$, the minimum is
$1-1/(2m)$, attained at $\alpha=1$.
\end{proposition}

\begin{proof}
The first affine function is strictly increasing and the second is
strictly decreasing. Equality between them is equivalent to
$\alpha(m+1)=m$. To the left of this intersection the second
function is larger, and to the right the first is larger.
Substitution proves \eqref{qr-eq-optimal-exponent} and uniqueness.
Since $m/(m+1)<1$, the first function is the larger throughout
$\alpha\geq1$; evaluating it at one proves the last assertion.
\end{proof}

Putting $H=D^\alpha$ in \eqref{qr-eq-power-extraction} gives, apart
from an arbitrarily small exponent loss, the bound $D^{E_m(\alpha)}$.
For sextic rows its best exponent over all positive row lengths would
therefore be $6/7$. The actually supplied range
$\alpha=1+\vartheta$, with $\vartheta>0$, has infimum $11/12$.
Thus changing this row length alone cannot obtain exponent $1/2$
from the two terms of \eqref{qr-eq-power-extraction}. This is a
limitation of that particular deduction. It does not disprove a
stronger arithmetic estimate, an alternative extraction, or GRH.
The stronger $7/8$ conclusion stated in \cite{OpenAISeptember}
uses additional estimates; it is not obtained by simply repeating
the extraction above.

\section{Completion and finite iteration}
\label{qr-sec-iteration}

\subsection{Removing added cube factors}

A completed theta expansion can contain additional cube factors.
Their removal is an exact algebraic step before any estimates are made.
The following formulation retains both shared prime factors and
possible zero values of the weights.

\begin{lemma}[M\"obius inversion for cube factors]\label{qr-lem-cube-inversion}
Let $w$ be a function on nonzero integral ideals such that
$w(\mathcal O_K)=1$ and $w(\mathfrak b\mathfrak c)
=w(\mathfrak b)w(\mathfrak c)$ for all ideals. Let $P:(0,\infty)
\to\mathbb C$ satisfy $P(X)=0$ for $0<X<X_0$, for some $X_0>0$.
Define
\begin{equation}\label{qr-eq-cube-completion}
 T(X)=\sum_{\mathfrak b}\frac{w(\mathfrak b)}{N\mathfrak b}
            P\left(\frac{X}{(N\mathfrak b)^3}\right).
\end{equation}
Then both this sum and the following sum are finite, and
\begin{equation}\label{qr-eq-cube-inversion}
 P(X)=\sum_{\mathfrak d}
       \frac{\mu_K(\mathfrak d)w(\mathfrak d)}{N\mathfrak d}
       T\left(\frac{X}{(N\mathfrak d)^3}\right).
\end{equation}
\end{lemma}

\begin{proof}
A nonzero summand in \eqref{qr-eq-cube-completion} must have
$(N\mathfrak b)^3\leq X/X_0$, so only finitely many ideals occur.
Also $T(X)=0$ for $X<X_0$, giving finiteness of the inverse sum.
Substitute \eqref{qr-eq-cube-completion} into its right-hand side
and group by $\mathfrak b=\mathfrak d\mathfrak c$. Complete
multiplicativity and multiplicativity of the norm make the coefficient
of $P(X/(N\mathfrak b)^3)$ equal to
\[
 \frac{w(\mathfrak b)}{N\mathfrak b}
       \sum_{\mathfrak d\mid\mathfrak b}\mu_K(\mathfrak d).
\]
For a nonunit ideal the last sum is the product of one factor
$1-1$ for each distinct prime ideal dividing it, and is zero.
For the unit ideal it is one. Since $w(\mathcal O_K)=1$, the
remaining term is exactly $P(X)$.
\end{proof}

No coprimality between $\mathfrak d$ and $\mathfrak c$ is assumed.
The proof still works when $w$ vanishes at some prime ideals.
In contrast, discarding the added cube terms from a complex sum is
not justified by positivity: they may cancel the other terms.
Splitting \eqref{qr-eq-cube-inversion} at a cutoff leaves a
remainder which must be estimated. This exact inversion is used in
\cite[Section~5]{OpenAIOctober}.
% Source labels: lem:cube-reduction, eq:cube-inverse.

\subsection{A finite descent with an explicit gap}

The next lemma isolates the elementary induction from any analytic
construction of successor sums. The requirement that every successor
remain in the same admissible class is part of its hypotheses.

\begin{lemma}\label{qr-lem-finite-descent}
Fix $B,a,\kappa>0$. For each $D\geq2$, let $\mathcal I(D)$ be a set
of indices. To each index $r$ associate numbers $Q_r\geq0$ and
$0<h_r\leq D^B$, and a set of successors
$\mathcal C(r)\subseteq\mathcal I(D)$. Suppose that
\[
 h_{r'}\leq D^{-a\kappa}h_r\qquad(r'\in\mathcal C(r)).
\]
Assume that for every $\delta>0$ there is $C_\delta\geq1$, independent
of $D$ and of the indices, such that
\begin{align*}
 Q_r&\leq C_\delta D^\delta &&\text{if }h_r\leq1,\\
 Q_r&\leq C_\delta D^\delta
           \left(1+\sup_{r'\in\mathcal C(r)}Q_{r'}\right)
       &&\text{if }h_r>1.
\end{align*}
An empty supremum is zero. Then for every $\varepsilon>0$ there is
$C^*_{\varepsilon,B,a,\kappa}$, independent of $D,r$, such that
$Q_r\leq C^*_{\varepsilon,B,a,\kappa}D^\varepsilon$.
\end{lemma}

\begin{proof}
Put $N=\lceil B/(a\kappa)\rceil$. After $N$ successor steps,
the row parameter is at most $D^{B-Na\kappa}\leq1$.
Fix $\varepsilon>0$, set $\delta=\varepsilon/(N+1)$, and write
$A=C_\delta D^\delta\geq1$. Starting with the terminal bound $A$
and applying the recurrence at most $N$ times gives
\[
 Q_r\leq\sum_{j=1}^{N+1}A^j
      \leq(N+1)C_\delta^{N+1}D^{(N+1)\delta}
      =(N+1)C_\delta^{N+1}D^\varepsilon.
\]
This argument is an induction on the maximum remaining number of
steps. It is valid even with infinitely many successors, because
the inductive bounds are uniform before each supremum is taken.
\end{proof}

In the October argument, the normalized quantity has the form
$Q=\mathcal E(\mathcal H,X,F)/(XF)$ and $\Sigma=XF$. The admissible
range includes $\mathcal H/\Sigma\leq D^{-\kappa}$. Its two analytic
inputs produce successors satisfying
\begin{equation}\label{qr-eq-source-descent}
 \mathcal H'<\mathcal H(\mathcal H/\Sigma)^2
       \leq D^{-2\kappa}\mathcal H,
 \qquad\frac{\mathcal H'}{\Sigma'}\leq
             \frac{\mathcal H}{\Sigma},
 \qquad\Sigma'\leq\Sigma.
\end{equation}
The second and third inequalities preserve these displayed size and gap
conditions, while the first provides the finite descent. Closure of the
other data under each transformation must be established separately. The analytic
proofs of those inputs remain separate obligations. The source
also tracks changing smooth tests: a derivative order $m$ is
replaced by an order at least $4m+12$, and the support interval
enlarges in a controlled way. For a fixed number of steps these
costs can be specified in advance. See
\cite[Section~5]{OpenAIOctober}.
% Source labels: prop:canonical, prop:transfer, sec:canonical-proof.

\begin{proposition}[Fixed parameters cannot be varied without bounds]
\label{qr-prop-nonuniform-kappa}
There exists a nonnegative function $F(D,\kappa)$ such that, for every
fixed $\kappa>0$ and every $\varepsilon>0$,
$F(D,\kappa)=O_{\kappa,\varepsilon}(D^\varepsilon)$ as $D\to\infty$,
but $F(D,1/\log D)=D$ for all $D\geq2$.
\end{proposition}

\begin{proof}
Take $F(D,\kappa)=\exp(1/\kappa)$. At fixed $\kappa$ it is a constant
in $D$, so the first estimate holds with constant $\exp(1/\kappa)$.
Substituting $\kappa=1/\log D$ proves the second assertion.
\end{proof}

Thus a statement proved for every fixed positive gap does not itself
control gaps tending to zero with $D$. In Lemma~\ref{qr-lem-finite-descent}
both the number of steps and the constant may deteriorate as
$\kappa\downarrow0$. To use a growing number of transformations one
must prove the needed quantitative uniformity, including all weight
and derivative costs. Proposition~\ref{qr-prop-nonuniform-kappa}
concerns that logical issue; it is not an assertion about the
actual arithmetic quantities of the source paper.

\section{Uniform continuation for a family of functions}
\label{qr-sec-uniform}

The September paper uses a different analytic arrangement: it
compares two estimates for a sum whose inverse Mellin representation
contains $1/L$. We record a precise abstract version, together with
the uniformity it requires. This is the continuation argument of
\cite[Section~2]{OpenAISeptember}, stated without its arithmetic
construction of the compared sums.
% Source labels: lem:continuation-criterion, eq:uniform-saving.

Let $(L_j)_{j\in\mathcal J}$ be a family of nonzero meromorphic
functions on $\mathbb C$. Suppose none has a zero in $\Re s>1$ and
that $1/L_j$ is bounded on every half-plane $\Re s\geq b>1$,
with the bound allowed to depend on $j,b$. Define
\[
 \beta_*:=\sup\left(\{1/2\}\cup
   \{\Re\rho:\Re\rho\geq1/2,\ L_j(\rho)=0
                        \text{ for some }j\}\right),
\]
where poles are not counted as zeros. Then $1/2\leq\beta_*\leq1$.

\begin{proposition}\label{qr-prop-uniform-continuation}
Fix $\sigma_0\in(1/2,1)$ and suppose $\beta_*>\sigma_0$.
Let $c,b_0\in\mathbb R$ be fixed. Suppose there are numbers
$\omega,\tau$ independent of $j$ with
$0<\omega<\beta_*-\sigma_0$ and $\tau>0$, and, for every $j$,
a holomorphic function $H_j$ on $\Re s>\sigma_0$ such that
$|H_j(s)-1|\leq1/2$ throughout that half-plane. Put
\begin{equation}\label{qr-eq-gaussian-mellin}
 f_j(Z)=\frac1{2\pi i}\int_{\Re s=2}
       Z^{s+c}e^{(s-b_0)^2}\frac{H_j(s)}{L_j(s)}\,ds
       \qquad(Z>0).
\end{equation}
Suppose functions $J_j(Z)$, defined for all sufficiently large $Z$,
satisfy
\begin{align}
 |J_j(Z)|&\ll_j Z^{\sigma_0+c+\omega},
       \label{qr-eq-first-comparison}\\
 |J_j(Z)-f_j(Z)|&\ll_j Z^{\beta_*+c-\tau}.
       \label{qr-eq-second-comparison}
\end{align}
The lower thresholds may depend on $j$. These hypotheses are
inconsistent with $\beta_*>\sigma_0$.
\end{proposition}

\begin{proof}
Let $\delta=\min(\beta_*-\sigma_0-\omega,\tau)>0$.
It is independent of $j$, and $\beta_*-\delta>\sigma_0$.
The triangle inequality gives
\begin{equation}\label{qr-eq-uniform-power-saving}
 |f_j(Z)|\ll_j Z^{\beta_*+c-\delta}
 \quad\text{for all sufficiently large }Z.
\end{equation}

We also need the other endpoint of its Mellin transform. The
integral \eqref{qr-eq-gaussian-mellin} converges absolutely, because
on every fixed vertical strip the Gaussian has modulus a bounded
factor times $e^{-(\Im s)^2}$. The reciprocal and $H_j$ are bounded
on $\Re s\geq2$. For any fixed $B>2$, Cauchy's theorem on rectangles
therefore moves the contour from $2$ to $B$. On the horizontal sides
the Gaussian dominates all remaining factors, so their integrals
tend to zero. Consequently
\begin{equation}\label{qr-eq-small-scale-decay}
 |f_j(Z)|\ll_{j,B}Z^{B+c}\qquad(0<Z\leq1).
\end{equation}
Together \eqref{qr-eq-uniform-power-saving} and
\eqref{qr-eq-small-scale-decay} show that
\[
 F_j(s)=\int_0^\infty f_j(Z)Z^{-s-c}\,\frac{dZ}{Z}
\]
is locally uniformly convergent and holomorphic in
$\Re s>\beta_*-\delta$. Near zero choose $B$ larger than the real
parts on the compact set under consideration; near infinity use
\eqref{qr-eq-uniform-power-saving}. Logarithmic factors introduced
by differentiation are integrable in the same way.

To identify the transform, set
$G_j(s)=e^{(s-b_0)^2}H_j(s)/L_j(s)$ for $\Re s>1$.
Writing $Z=e^u$ in \eqref{qr-eq-gaussian-mellin} gives
\[
 e^{-(2+c)u}f_j(e^u)
   =\frac1{2\pi}\int_{\mathbb R}e^{itu}G_j(2+it)\,dt.
\]
The integrand $G_j(2+it)$ is continuous and integrable. The
left-hand side is integrable in $u$: at $+\infty$ use
$\beta_*\leq1$ and \eqref{qr-eq-uniform-power-saving}, and at
$-\infty$ use \eqref{qr-eq-small-scale-decay} with $B>2$.
Fourier inversion therefore gives
$F_j(2+it)=G_j(2+it)$ for every real $t$.
Both functions are holomorphic on $\Re s>1$, so the identity
theorem gives $F_j(s)=G_j(s)$ throughout that half-plane.

Since $|H_j|\geq1/2$, the function
\[
 R_j(s):=e^{-(s-b_0)^2}\frac{F_j(s)}{H_j(s)}
\]
is a holomorphic continuation of $1/L_j(s)$ to
$\Re s>\beta_*-\delta$. The same $\delta>0$ works for every $j$.
By the definition of the supremum and
$\beta_*-\delta>\sigma_0>1/2$, some $L_j$ has a zero $\rho$ with
$\Re\rho>\beta_*-\delta$. Its reciprocal has a pole there, in
contradiction to this holomorphic continuation. Equivalently, the
meromorphic identity theorem gives $L_j(s)R_j(s)=1$ throughout
the larger half-plane; evaluating at $\rho$ gives $0=1$.
\end{proof}

The constants in the two comparisons need not be uniform in $j$.
The positive saving $\delta$ does need to be uniform. The next
example proves that this distinction cannot be omitted.

\begin{proposition}\label{qr-prop-nonuniform-family}
There is a family of entire functions with zero real parts having
supremum one, although the reciprocal of each member extends
holomorphically from $\Re s>1$ to a strictly larger half-plane.
There is no common positive amount of extension for the whole family.
\end{proposition}

\begin{proof}
For integers $n\geq1$, let
\[
 \rho_n=1-\frac1{n+2},\qquad L_n(s)=s-\rho_n,
 \qquad\delta_n=\frac1{2(n+2)}.
\]
The only zero of $L_n$ is $\rho_n=1-2\delta_n$.
Hence $1/L_n(s)$ is holomorphic on $\Re s>1-\delta_n$,
which extends a positive distance to the left of one for each
fixed $n$. Nevertheless $\rho_n\to1$, so the supremum of their
real parts is one. Given any $\delta>0$, choose $n$ with
$1/(n+2)<\delta$; then $\rho_n>1-\delta$, so the $n$th reciprocal
cannot be holomorphic on $\Re s>1-\delta$.
\end{proof}

These are elementary entire functions, not arithmetic $L$-functions.
Their purpose is to exhibit the quantifier error in replacing a
common positive saving by separate savings which may tend to zero.
In an arithmetic application, the uniform saving and the two
comparisons in Proposition~\ref{qr-prop-uniform-continuation} must
be established by the actual estimates. The continuation argument
does not improve a fixed $11/12$ or $7/8$ half-plane to $1/2$
without those new estimates.
